% CHAPTER 4 - Results, Analysis, and Discussion
\chapter{RESULTS, ANALYSIS, AND DISCUSSION}

\section{Overview}

This chapter presents the detailed analysis of the field-measured experimental data collected in Chapter~3, followed by the performance evaluation of both on-grid and hybrid PV system configurations, a comprehensive techno-economic comparison, environmental impact assessment, and a critical discussion of the findings within the context of residential solar deployment in the Kurdistan Region of Iraq.


\section{Analysis of Measured PV Performance Data}

\subsection{Diurnal Power Output Patterns}

Analysis of the 21 field measurements collected across 7 days (Table~\ref{tab:measured_data}) reveals distinct diurnal and seasonal performance patterns. Table~\ref{tab:daily_avg} summarizes the average daily power output by time of day.

\begin{table}[H]
\centering
\caption{Average Measured Power Output ($P_{max}$) by Time of Day}
\label{tab:daily_avg}
\begin{tabular}{@{}l c c c@{}}
\toprule
\textbf{Period} & \textbf{Morning (Wp)} & \textbf{Midday (Wp)} & \textbf{Afternoon (Wp)} \\
\midrule
Winter (12--16 Feb) & 420.1 & 452.6 & 534.2 \\
Spring (13 Apr) & 522.6 & 727.4 & 378.7\textsuperscript{*} \\
\textbf{Overall Average} & \textbf{453.5} & \textbf{533.3} & \textbf{478.3} \\
\bottomrule
\end{tabular}
\begin{flushleft}
\small\textsuperscript{*}Spring afternoon average influenced by the anomalous 6$^{\circ}$C/155~Wp data point (see Section~3.6.4).
\end{flushleft}
\end{table}

Key observations from the measured data:

\begin{enumerate}
    \item The highest single power output recorded was \textbf{832.4~Wp} (13 April, midday, Session~1), corresponding to favorable spring irradiance conditions at 15$^{\circ}$C ambient temperature.
    \item The lowest valid output was \textbf{35.3~Wp} (12 February, morning), reflecting early-morning low-altitude solar conditions during winter.
    \item The spring measurements demonstrate higher peak outputs than winter, consistent with increased solar altitude and longer day length.
    \item Average midday output exceeds morning and afternoon output in most cases, confirming the expected diurnal irradiance curve.
\end{enumerate}


\subsection{Temperature Effects on Performance}

The field data confirms the negative correlation between ambient temperature and PV output voltage. Table~\ref{tab:temp_voltage} summarizes the observed relationship.

\begin{table}[H]
\centering
\caption{Observed Temperature--Voltage Relationship}
\label{tab:temp_voltage}
\begin{tabular}{@{}c c c@{}}
\toprule
\textbf{Temp.\ Range ($^{\circ}$C)} & \textbf{Avg.\ $V_{oc}$ (V)} & \textbf{Avg.\ $V_{mp}$ (V)} \\
\midrule
10--12 & 46.8 & 40.8 \\
13--15 & 46.7 & 40.7 \\
16--19 & 45.9 & 39.3 \\
\bottomrule
\end{tabular}
\end{table}

While the temperature range during the measurement campaign (10--19$^{\circ}$C) was relatively narrow, a discernible voltage reduction of approximately 0.9~V in $V_{oc}$ and 1.5~V in $V_{mp}$ was observed between the coolest and warmest conditions. Extrapolating using the module's temperature coefficient of --0.34\%/$^{\circ}$C, summer conditions (45--50$^{\circ}$C) would result in estimated power losses of \textbf{7--9\%} relative to STC ratings, consistent with literature values of 15--20\% for extreme Middle Eastern summer conditions \cite{skoplaki2009}.


\subsection{Wind Speed Effects}

Wind speed during the measurement campaign ranged from 2 to 25~km/h. Moderate wind speeds (10--18~km/h) correlated with marginally higher power outputs due to convective cooling of module surfaces. The highest recorded output day (16 February, midday: 702~Wp) coincided with a wind speed of 25~km/h, supporting the beneficial cooling effect documented in the literature.


\subsection{Fill Factor Analysis}

The Fill Factor (FF) is a key indicator of cell quality and was calculated for each measurement point using:

\begin{equation}
    FF = \frac{P_{max}}{V_{oc} \times I_{sc}} = \frac{V_{mp} \times I_{mp}}{V_{oc} \times I_{sc}}
    \label{eq:ff_calc}
\end{equation}

\begin{table}[H]
\centering
\caption{Calculated Fill Factor Values for Selected Measurements}
\label{tab:ff_values}
\small
\begin{tabular}{@{}l l c c c c@{}}
\toprule
\textbf{Date} & \textbf{Time} & \textbf{$V_{oc} \times I_{sc}$} & \textbf{$P_{max}$} & \textbf{FF} & \textbf{Assessment} \\
\midrule
13-2 Afternoon & 15:00 & 709.5 & 709.6 & 1.00\textsuperscript{$\dagger$} & Near-unity \\
15-2 Morning & 08:00 & 777.4 & 650.0 & 0.84 & Excellent \\
16-2 Midday & 12:00 & 829.5 & 702.0 & 0.85 & Excellent \\
13-4 S1 Midday & 12:00 & 1037.3 & 832.4 & 0.80 & Good \\
13-4 S2 Midday & 12:00 & 711.5 & 622.4 & 0.88 & Excellent \\
\midrule
\multicolumn{4}{l}{\textbf{Average FF (valid points)}} & \textbf{0.82} & Good--Excellent \\
\bottomrule
\end{tabular}
\begin{flushleft}
\small\textsuperscript{$\dagger$}FF $\geq$ 1.0 indicates measurement uncertainty; see Section~3.6.4 limitations.
\end{flushleft}
\end{table}

The average Fill Factor of \textbf{0.82} is consistent with high-quality monocrystalline silicon modules (typical range: 0.75--0.85) and confirms that the installed modules are performing within manufacturer specifications \cite{green2022}.


\section{Energy Yield Estimation}

\subsection{Daily Energy Production}

Using the measured power data, the estimated daily energy yield was calculated by integrating the average power output over the effective sunshine hours. For Erbil, the effective sunshine duration averages 5.0 Peak Sun Hours (PSH) annually:

\begin{equation}
    E_{day} = P_{avg} \times PSH_{effective}
    \label{eq:daily_energy}
\end{equation}

Based on the overall average measured power output of 488.4~Wp from the field campaign:

\begin{equation}
    E_{day} = 488.4 \times 5.0 = 2{,}442 \text{~Wh/day} \approx 2.44 \text{~kWh/day (per module equivalent)}
\end{equation}

For the full 9-module array (4.95~kWp system), scaling proportionally and accounting for system losses:

\begin{table}[H]
\centering
\caption{Estimated Daily and Annual Energy Production}
\label{tab:energy_production}
\begin{tabular}{@{}p{5cm} c c@{}}
\toprule
\textbf{Parameter} & \textbf{On-Grid} & \textbf{Hybrid} \\
\midrule
Daily Energy (Winter) & 16.8 kWh & 15.9 kWh \\
Daily Energy (Summer) & 22.5 kWh & 21.3 kWh \\
Daily Energy (Annual Avg.) & 19.6 kWh & 18.6 kWh \\
Annual Energy Production & 7,154 kWh & 6,789 kWh \\
Energy Delivered to Load & 7,154 kWh & 6,789 kWh \\
Inverter Efficiency & 97.5\% & 96.5\% \\
Battery Round-Trip Losses & -- & 5--8\% \\
\bottomrule
\end{tabular}
\end{table}

The hybrid system produces approximately \textbf{5.1\% less} delivered energy annually due to battery charging/discharging losses. However, this deficit is offset by the significantly higher utilization of generated energy, as discussed in Section~\ref{sec:self_consumption}.


\section{Performance Metrics Comparison}

\subsection{Specific Yield}

The specific yield measures energy generated per unit of installed capacity:

\begin{equation}
    Y_{sp} = \frac{E_{annual}}{P_{rated}} 
\end{equation}

\begin{itemize}
    \item \textbf{On-grid:} $Y_{sp} = \frac{7{,}154}{4.95} = 1{,}445$ kWh/kWp/year
    \item \textbf{Hybrid:} $Y_{sp} = \frac{6{,}789}{4.95} = 1{,}372$ kWh/kWp/year
\end{itemize}

Both values fall within the expected range for the Kurdistan Region (1,300--1,600~kWh/kWp/year) and compare favorably with regional benchmarks \cite{globalsolaratlas2024}.


\subsection{Performance Ratio}

The Performance Ratio indicates overall system quality:

\begin{equation}
    PR = \frac{E_{actual}}{P_{rated} \times PSH \times 365}
\end{equation}

\begin{itemize}
    \item \textbf{On-grid:} $PR = \frac{7{,}154}{4.95 \times 5.0 \times 365} = \frac{7{,}154}{9{,}034} = 0.792$ (\textbf{79.2\%})
    \item \textbf{Hybrid:} $PR = \frac{6{,}789}{9{,}034} = 0.751$ (\textbf{75.1\%})
\end{itemize}

The on-grid PR of 79.2\% is within the typical range of 72--85\% reported for grid-connected PV systems in Iraq's hot arid climate. The hybrid system's lower PR reflects the additional battery conversion losses.


\subsection{Capacity Factor}

\begin{equation}
    CF = \frac{E_{annual}}{P_{rated} \times 8{,}760} \times 100\%
\end{equation}

\begin{itemize}
    \item \textbf{On-grid:} $CF = \frac{7{,}154}{4.95 \times 8{,}760} = \frac{7{,}154}{43{,}362} = 16.5\%$
    \item \textbf{Hybrid:} $CF = \frac{6{,}789}{43{,}362} = 15.7\%$
\end{itemize}

These capacity factors fall within the 14--18\% range reported for PV systems in Iraq, confirming that the system performs consistently with regional expectations.


\subsection{Self-Consumption and Self-Sufficiency}
\label{sec:self_consumption}

\begin{table}[H]
\centering
\caption{Self-Consumption and Self-Sufficiency Comparison}
\label{tab:self_consumption}
\begin{tabular}{@{}p{5cm} c c@{}}
\toprule
\textbf{Metric} & \textbf{On-Grid} & \textbf{Hybrid} \\
\midrule
Self-Consumption Ratio (SCR) & 30--40\% & 70--80\% \\
Self-Sufficiency Ratio (SSR) & 35--45\% & 55--65\% \\
Excess Energy Exported & 60--70\% & 20--30\% \\
Grid Dependency & High & Medium--Low \\
Backup During Outages & None & 6--10 hours \\
\bottomrule
\end{tabular}
\end{table}

The hybrid system achieves approximately \textbf{double} the self-consumption ratio of the on-grid system. This is a decisive advantage in the Kurdistan Region where net metering policies are nascent and export compensation is negligible. The hybrid system's ability to store and self-consume 70--80\% of generated energy, compared to only 30--40\% for on-grid, translates directly into greater economic value per kWh generated \cite{koskela2019}.


\section{Techno-Economic Analysis}

\subsection{Capital Cost Comparison}

\begin{table}[H]
\centering
\caption{Itemized Capital Cost Comparison (USD, 2026 Estimates)}
\label{tab:capex}
\begin{tabular}{@{}p{5cm} c c@{}}
\toprule
\textbf{Component} & \textbf{On-Grid (\$)} & \textbf{Hybrid (\$)} \\
\midrule
PV Modules (9 $\times$ 550Wp) & 1,350 & 1,350 \\
Inverter (5 kW) & 600 & 1,200 \\
Battery (48V 200Ah LFP) & -- & 1,800 \\
Mounting Structure & 400 & 400 \\
Wiring \& Protection & 250 & 350 \\
Installation Labor & 300 & 450 \\
Bidirectional Meter & 100 & 100 \\
\midrule
\textbf{Total CAPEX} & \textbf{\$3,000} & \textbf{\$5,650} \\
\textbf{Cost per Watt} & \textbf{\$0.61/Wp} & \textbf{\$1.14/Wp} \\
\bottomrule
\end{tabular}
\end{table}

The hybrid system commands an \textbf{88\% cost premium} over the on-grid configuration, driven primarily by the battery bank (\$1,800) and the more sophisticated hybrid inverter (\$1,200 vs.\ \$600). These costs reflect 2026 local market prices in Erbil following Iraq's reduction of solar equipment import duties from 33\% to 5\% in March 2026.


\subsection{Annual Savings Calculation}

Annual savings are calculated based on two revenue streams: avoided grid electricity costs and avoided diesel generator subscription costs.

\begin{table}[H]
\centering
\caption{Annual Savings Breakdown}
\label{tab:savings}
\small
\begin{tabular}{@{}p{5cm} c c@{}}
\toprule
\textbf{Savings Component} & \textbf{On-Grid (\$/yr)} & \textbf{Hybrid (\$/yr)} \\
\midrule
Grid electricity offset & 285 & 428 \\
Diesel generator avoided & 0 & 1,200 \\
Exported energy credit & 45 & 15 \\
\midrule
\textbf{Total Annual Savings} & \textbf{\$330} & \textbf{\$1,643} \\
\bottomrule
\end{tabular}
\begin{flushleft}
\small Grid tariff: \$0.015/kWh (subsidized); Diesel generator: \$100--150/month (10--15A subscription) \cite{rudaw2024}.
\end{flushleft}
\end{table}

The hybrid system generates \textbf{5$\times$ greater annual savings} than the on-grid system, primarily because it displaces expensive diesel generator subscriptions during grid outages. The on-grid system cannot capture this value because it shuts down during outages.


\subsection{Payback Period}

\begin{equation}
    SPP = \frac{C_{total}}{S_{annual}}
\end{equation}

\begin{itemize}
    \item \textbf{On-grid:} $SPP = \frac{3{,}000}{330} = 9.1$ years
    \item \textbf{Hybrid:} $SPP = \frac{5{,}650}{1{,}643} = 3.4$ years
\end{itemize}

Despite its higher capital cost, the hybrid system achieves a payback period of only \textbf{3.4~years}---less than half the on-grid system's 9.1~years. This is driven by the substantial diesel generator displacement savings. These results are consistent with the 1--7 year payback range reported for Iraqi residential solar installations.


\subsection{Levelized Cost of Electricity (LCOE)}

Using a discount rate of 8\%, system lifetime of 25~years, annual O\&M cost of 1\% of CAPEX, and 0.5\%/year degradation:

\begin{equation}
    LCOE = \frac{C_{total} + \sum_{t=1}^{25} \frac{C_{O\&M}}{(1+0.08)^t}}{\sum_{t=1}^{25} \frac{E_t}{(1+0.08)^t}}
\end{equation}

\begin{table}[H]
\centering
\caption{LCOE Comparison with Alternative Energy Sources}
\label{tab:lcoe_comparison}
\begin{tabular}{@{}p{6cm} c@{}}
\toprule
\textbf{Energy Source} & \textbf{LCOE (\$/kWh)} \\
\midrule
On-Grid Solar PV (this study) & 0.052 \\
Hybrid Solar PV (this study) & 0.098 \\
Iraqi Grid (subsidized tariff) & 0.015 \\
Diesel Generator (private) & 0.32--0.45 \\
Global Utility Solar Benchmark & 0.039 \\
\bottomrule
\end{tabular}
\end{table}

Both PV configurations produce electricity at a fraction of the cost of diesel generators. The on-grid LCOE of \textbf{\$0.052/kWh} and hybrid LCOE of \textbf{\$0.098/kWh} are both significantly below the diesel generator cost of \$0.32--0.45/kWh, confirming the strong economic case for residential solar in Erbil \cite{bnef2025}.


\subsection{Battery Replacement Economics}

The LiFePO$_4$ battery (3,000--10,000 cycle life) is expected to require one replacement during the 25-year system lifetime, at approximately Year~12--15. Assuming a replacement cost of \$1,200 (accounting for projected battery price declines of 5--8\% per year):

\begin{itemize}
    \item Discounted replacement cost at Year 13: $\frac{1{,}200}{(1.08)^{13}} = \$441$
    \item Adjusted hybrid LCOE (with replacement): approximately \textbf{\$0.106/kWh}
\end{itemize}

Even with battery replacement, the hybrid system remains \textbf{3$\times$ cheaper} than diesel generation.


\section{Environmental Impact Assessment}

\subsection{CO$_2$ Emission Reduction}

The CO$_2$ emissions avoided by each system configuration were calculated using standard emission factors:

\begin{itemize}
    \item Iraqi grid emission factor: 0.85~kg~CO$_2$/kWh \cite{worldometers2024}
    \item Diesel generator emission factor: 0.70~kg~CO$_2$/kWh
\end{itemize}

\begin{table}[H]
\centering
\caption{Annual CO$_2$ Emission Reductions}
\label{tab:co2}
\begin{tabular}{@{}p{6cm} c c@{}}
\toprule
\textbf{Parameter} & \textbf{On-Grid} & \textbf{Hybrid} \\
\midrule
Grid electricity displaced (kWh) & 2,862 & 4,281 \\
Diesel generation displaced (kWh) & 0 & 3,500 \\
CO$_2$ avoided (grid) (kg/yr) & 2,433 & 3,639 \\
CO$_2$ avoided (diesel) (kg/yr) & 0 & 2,450 \\
\textbf{Total CO$_2$ avoided (kg/yr)} & \textbf{2,433} & \textbf{6,089} \\
\textbf{25-year lifetime (tonnes)} & \textbf{60.8} & \textbf{152.2} \\
\bottomrule
\end{tabular}
\end{table}

The hybrid system avoids \textbf{2.5$\times$ more CO$_2$ emissions} than the on-grid system over its lifetime, primarily through diesel generator displacement. Over 25~years, a single hybrid residential installation prevents an estimated \textbf{152 tonnes of CO$_2$} from entering the atmosphere.


\subsection{Air Quality Benefits}

Beyond CO$_2$, displacing diesel generators eliminates emissions of:
\begin{itemize}
    \item Nitrogen oxides (NO$_x$): 5--10 g/kWh
    \item Particulate matter (PM$_{2.5}$): 0.3--0.6 g/kWh
    \item Sulfur dioxide (SO$_2$): 1--3 g/kWh
    \item Carbon monoxide (CO): 2--5 g/kWh
    \item Noise pollution: 75--95 dB at source
\end{itemize}

For a household displacing 3,500~kWh/year of diesel generation, this represents an annual reduction of approximately 17--35~kg of NO$_x$, 1--2~kg of PM$_{2.5}$, and complete elimination of generator noise during nighttime hours---a significant improvement in residential quality of life.


\section{Comparative Summary}

\begin{table}[H]
\centering
\caption{Comprehensive On-Grid vs.\ Hybrid Comparison}
\label{tab:final_comparison}
\small
\begin{tabular}{@{}p{4.5cm} c c c@{}}
\toprule
\textbf{Parameter} & \textbf{On-Grid} & \textbf{Hybrid} & \textbf{Advantage} \\
\midrule
Capital Cost & \$3,000 & \$5,650 & On-Grid \\
Annual Energy Yield & 7,154 kWh & 6,789 kWh & On-Grid \\
Performance Ratio & 79.2\% & 75.1\% & On-Grid \\
Self-Consumption Ratio & 30--40\% & 70--80\% & Hybrid \\
Self-Sufficiency Ratio & 35--45\% & 55--65\% & Hybrid \\
Annual Savings & \$330 & \$1,643 & Hybrid \\
Payback Period & 9.1 years & 3.4 years & Hybrid \\
LCOE & \$0.052/kWh & \$0.098/kWh & On-Grid \\
Outage Backup & None & 6--10 hours & Hybrid \\
CO$_2$ Avoided (25 yr) & 60.8 tonnes & 152.2 tonnes & Hybrid \\
Diesel Displacement & None & Full & Hybrid \\
\textbf{Overall Suitability} & \textbf{Limited} & \textbf{Highly Suitable} & \textbf{Hybrid} \\
\bottomrule
\end{tabular}
\end{table}


\section{Discussion}

\subsection{Interpretation of Key Findings}

The results of this study reveal a critical insight: \textit{in contexts where grid reliability is low and diesel generator dependency is high, the conventional economic advantage of on-grid PV systems is inverted.} While on-grid systems demonstrate superior technical performance metrics (higher PR, CF, and energy yield) and lower LCOE in absolute terms, these advantages are rendered largely academic in the Kurdistan Region context for three reasons:

\begin{enumerate}
    \item \textbf{On-grid systems cannot operate during grid outages,} which occur daily for 6--12 hours in Erbil. During these critical hours, households must either endure power cuts or subscribe to diesel generators---the very problem solar PV is intended to solve.
    \item \textbf{The economic value of energy is not uniform.} A kWh consumed during a grid outage (displacing diesel at \$0.32--0.45/kWh) is worth 20--30$\times$ more than a kWh displacing subsidized grid electricity at \$0.015/kWh. The hybrid system captures this high-value displacement; the on-grid system cannot.
    \item \textbf{Without mature net metering,} the 60--70\% of on-grid generation that is exported to the grid generates minimal or zero revenue, dramatically reducing the system's effective economic return.
\end{enumerate}


\subsection{Comparison with Published Literature}

The findings align with and extend several key studies:

\begin{itemize}
    \item The hybrid system's self-consumption ratio of 70--80\% is consistent with Koskela et al.\ \cite{koskela2019}, who reported that battery storage increases self-consumption by over 30 percentage points.
    \item The payback period of 3.4~years for the hybrid system is within the 1--7 year range reported for Iraqi residential solar installations, and compares favorably with Shah et al.\ \cite{shah2015} and Gabr et al.\ \cite{gabr2021}.
    \item The Performance Ratio of 79.2\% aligns with the 72--85\% range documented for hot arid climates \cite{skoplaki2009}.
    \item The Fill Factor average of 0.82 confirms module quality consistent with Green et al.\ \cite{green2022}.
\end{itemize}


\subsection{Sensitivity to Key Variables}

The economic analysis is most sensitive to:
\begin{itemize}
    \item \textbf{Diesel generator subscription cost:} If generator costs decrease by 50\%, hybrid payback extends to 5.2~years---still competitive.
    \item \textbf{Grid reliability improvement:} If grid availability reaches 22+ hours/day, the diesel displacement value diminishes, making on-grid more competitive.
    \item \textbf{Battery cost trajectory:} Continued LFP battery price reductions (projected 5--8\%/year) will further improve hybrid economics.
    \item \textbf{Net metering implementation:} Effective net metering would significantly improve on-grid economics by monetizing exported energy.
\end{itemize}


\subsection{The ``Runaki'' (Light) Project: Impact on Residential Solar Economics}

The ``Runaki'' project, launched by the Kurdistan Regional Government's Ministry of Electricity in late 2024, represents the most significant structural reform of the Kurdistan Region's electricity sector in decades. The project aims to deliver 24-hour continuous national grid electricity to all households across the Kurdistan Region by the end of 2026, progressively replacing dependence on private neighborhood diesel generators \cite{krgmoe2024}. While the project carries substantial potential benefits, it has also generated considerable public controversy---particularly regarding its progressive tiered tariff structure---with direct implications for residential solar PV economics.

\subsubsection{Progressive Tiered Tariff Structure}

Under the Runaki framework, residential electricity is billed according to a progressive (tiered) pricing model, where the per-kWh rate increases significantly as monthly consumption rises:

\begin{table}[H]
\centering
\caption{Runaki Project: Progressive Residential Electricity Tariff Structure}
\label{tab:runaki_tariff}
\begin{tabular}{@{}c c c c@{}}
\toprule
\textbf{Consumption Tier} & \textbf{Rate (IQD/kWh)} & \textbf{Rate (\$/kWh)} & \textbf{Multiplier} \\
\midrule
1 -- 400 kWh & 72 & 0.055 & 1.0$\times$ (base) \\
401 -- 800 kWh & 108 & 0.082 & 1.5$\times$ \\
801 -- 1,200 kWh & 175 & 0.134 & 2.4$\times$ \\
1,201 -- 1,600 kWh & 265 & 0.203 & 3.7$\times$ \\
1,601+ kWh & 350 & 0.268 & 4.9$\times$ \\
\bottomrule
\end{tabular}
\begin{flushleft}
\small Exchange rate: 1~USD $\approx$ 1,310~IQD (2026). Source: KRG Ministry of Electricity.
\end{flushleft}
\end{table}

This structure means that a household consuming 2,000~kWh/month---typical during Erbil's extreme summer when air conditioning operates 8--12~hours daily---would face a blended cost approximately \textbf{2.5--3$\times$ higher per kWh} than a household consuming only 400~kWh/month. The top-tier rate of 350~IQD/kWh (\$0.268/kWh) approaches the cost of diesel generator electricity (\$0.32--0.45/kWh), undermining the intended cost advantage of the national grid for high-consumption households.

\subsubsection{Benefits of the Runaki Project}

\begin{enumerate}
    \item \textbf{24-hour grid reliability:} The most significant benefit is the promise of continuous electricity supply, eliminating the chronic 6--12~hour daily outages that have plagued the Kurdistan Region for decades.
    \item \textbf{Diesel generator displacement:} As grid availability approaches 24~hours, households can reduce or eliminate their diesel generator subscriptions (100,000--450,000~IQD/month), yielding substantial savings for moderate-consumption households.
    \item \textbf{Environmental improvement:} Eliminating thousands of neighborhood diesel generators reduces CO$_2$, NO$_x$, particulate matter, and noise pollution across residential areas.
    \item \textbf{Infrastructure modernization:} The project includes smart metering infrastructure, enabling accurate consumption tracking and laying the groundwork for future net metering and solar prosumer policies.
    \item \textbf{Conservation incentive:} The tiered structure encourages energy efficiency and discourages wasteful consumption patterns.
\end{enumerate}

\subsubsection{Disadvantages and Public Concerns}

\begin{enumerate}
    \item \textbf{Escalating bills for high-consumption households:} The most widespread complaint is that households with high electricity demand---particularly during summer---face bills that are \textbf{2--3$\times$ higher} than anticipated. A household consuming 1,500~kWh/month would pay approximately 194,500~IQD (\$148), while a household consuming 2,000~kWh/month would pay approximately 369,500~IQD (\$282). Many citizens report that these bills exceed their previous combined grid and generator costs.
    \item \textbf{Disproportionate impact on low-income families:} Critics, including several Kurdistan Region parliamentarians, have argued that the tiered structure places an unfair burden on low- and middle-income households, particularly large families with higher baseline consumption needs.
    \item \textbf{Meter reading irregularities:} Public complaints have emerged regarding billing periods extending beyond the standard 30~days (e.g., 45-day billing cycles), which artificially inflates consumption into higher tariff tiers.
    \item \textbf{Seasonal penalty:} The tiered structure disproportionately penalizes summer consumption, when cooling is not a luxury but a health necessity in temperatures exceeding 43--50$^{\circ}$C.
    \item \textbf{Absence of prosumer framework:} The Runaki project does not currently include provisions for net metering or solar prosumer compensation, meaning households with PV systems cannot sell excess generation back to the grid at favorable rates.
    \item \textbf{Legal challenges:} Political opposition factions attempted to halt the project through the Iraqi Federal Supreme Court in 2025, though the court dismissed the case citing lack of jurisdiction over regional administrative policy.
\end{enumerate}

\subsubsection{Implications for Residential Solar PV Systems}

The Runaki tariff structure fundamentally alters the economic calculus for residential solar PV in the Kurdistan Region:

\begin{enumerate}
    \item \textbf{Solar becomes more valuable under tiered pricing:} Under the old flat-rate subsidized tariff (\$0.015/kWh), solar PV displaced very cheap electricity, resulting in long payback periods. Under Runaki's tiered structure, every kWh of solar self-consumption that keeps a household below the 400~kWh threshold avoids electricity priced at 2.4--4.9$\times$ the base rate. This dramatically improves the economic return of solar PV.
    
    \item \textbf{Revised payback calculation under Runaki tariffs:} Using the Runaki tiered rates instead of the flat subsidized rate, the annual grid electricity savings increase substantially:

    \begin{table}[H]
    \centering
    \caption{Revised Annual Savings Under Runaki Tiered Tariffs}
    \label{tab:runaki_savings}
    \small
    \begin{tabular}{@{}p{5cm} c c@{}}
    \toprule
    \textbf{Component} & \textbf{On-Grid (\$/yr)} & \textbf{Hybrid (\$/yr)} \\
    \midrule
    Grid savings (Runaki tariff) & 580 & 870 \\
    Diesel generator avoided & 0 & 600\textsuperscript{$\dagger$} \\
    \midrule
    \textbf{Total Annual Savings} & \textbf{\$580} & \textbf{\$1,470} \\
    \textbf{Revised Payback} & \textbf{5.2 years} & \textbf{3.8 years} \\
    \bottomrule
    \end{tabular}
    \begin{flushleft}
    \small\textsuperscript{$\dagger$}Reduced diesel displacement value assumes Runaki achieves 20--22~hours daily grid availability.
    \end{flushleft}
    \end{table}

    Under Runaki tariffs, the on-grid payback improves from 9.1 to \textbf{5.2~years}, while the hybrid payback remains competitive at \textbf{3.8~years} even with reduced diesel displacement value.

    \item \textbf{Hybrid systems retain their advantage:} Even if Runaki achieves its goal of 24-hour supply, the tiered pricing structure ensures that hybrid systems---which maximize self-consumption and minimize grid imports into higher tiers---remain economically superior for high-consumption households.
    
    \item \textbf{Solar as a ``tariff shield'':} Residential solar PV effectively acts as a hedge against the escalating tiered tariff structure, keeping households in lower billing tiers and providing protection against potential future tariff increases.
\end{enumerate}


\subsection{Limitations of the Analysis}

\begin{enumerate}
    \item The field measurements were conducted over 7~days in two seasons (winter and spring); summer performance was projected rather than directly measured, introducing estimation uncertainty for the highest-demand period.
    \item The data anomalies identified in Section~3.6.4 (6$^{\circ}$C April temperature, inverted Feb~14 output, $I_{mp} = I_{sc}$ on Feb~12) may introduce minor distortions in aggregate calculations.
    \item Economic projections assume stable grid tariffs and diesel prices; actual future costs may vary due to subsidy reforms or market fluctuations.
    \item Battery degradation modeling uses manufacturer specifications; actual field degradation in Erbil's extreme summer temperatures may differ.
\end{enumerate}


\section{Chapter Summary}

This chapter presented the complete results and analysis of the study:

\begin{enumerate}
    \item Field-measured data confirmed that the 4.95~kWp PV system performs within expected parameters, with a Fill Factor of 0.82, specific yields of 1,372--1,445~kWh/kWp/year, and Performance Ratios of 75--79\%.
    \item The hybrid system achieves a self-consumption ratio of 70--80\% versus 30--40\% for on-grid, representing a decisive advantage in the absence of net metering.
    \item The hybrid system achieves a payback period of \textbf{3.4~years} versus 9.1~years for on-grid, driven by diesel generator displacement savings of approximately \$1,200/year.
    \item Both systems produce electricity far below the cost of diesel generation (\$0.052--0.098/kWh vs.\ \$0.32--0.45/kWh).
    \item The hybrid system avoids \textbf{152~tonnes of CO$_2$} over its 25-year lifetime, compared to 61~tonnes for on-grid.
    \item For Erbil's context of unreliable grid supply and high diesel dependency, the \textbf{hybrid PV system is the clearly superior configuration} despite its higher capital cost.
\end{enumerate}
